% microtype with pdflatex (PDF mode): font expansion (\pdffontexpand
% autoexpand, \pdfadjustspacing=2) and character protrusion
% (\pdfprotrudechars=2), microtype's defaults; for the pdflatex harness
% (tests/pdflatex_docs.rs): the PDF and the written files as pdflatex
% makes them, at every schedule (CORE.md, "Font expansion and character
% protrusion").
\documentclass{article}
\usepackage[T1]{fontenc}
\usepackage{lmodern}
\usepackage{microtype}
\begin{document}
\section{Expansion and protrusion}

A quick brown fox jumps over the lazy dog, and then (again) it jumps over
the lazy dog once more. ``Quick,'' said the fox --- and it jumped. The end
of the story: dogs are lazy, foxes are quick; and so on, and so forth,
until the end of the paragraph is reached at last, which it is, here.

\begin{minipage}{0.4\textwidth}
A narrow minipage stretches its fonts more: ``margins,'' hyphens-, and dots.
More text here, to have more lines; and still more, until it ends here.
\end{minipage}

\microtypesetup{expansion=false}
The same words without expansion: ``margins,'' hyphens-, and dots, as the
lines without it set them, to compare; and so on until the paragraph ends.
\end{document}
